%%% Hlavní soubor. Zde se definují základní parametry a odkazuje se na ostatní části. %%%

%% Verze pro jednostranný tisk:
% Okraje: levý 40mm, pravý 25mm, horní a dolní 25mm
% (ale pozor, LaTeX si sám přidává 1in)
\documentclass[12pt,a4paper]{report}
\setlength\textwidth{145mm}
\setlength\textheight{247mm}
\setlength\oddsidemargin{15mm}
\setlength\evensidemargin{15mm}
\setlength\topmargin{0mm}
\setlength\headsep{0mm}
\setlength\headheight{0mm}
% \openright zařídí, aby následující text začínal na pravé straně knihy
\let\openright=\clearpage

%% Pokud tiskneme oboustranně:
% \documentclass[12pt,a4paper,twoside,openright]{report}
% \setlength\textwidth{145mm}
% \setlength\textheight{247mm}
% \setlength\oddsidemargin{15mm}
% \setlength\evensidemargin{0mm}
% \setlength\topmargin{0mm}
% \setlength\headsep{0mm}
% \setlength\headheight{0mm}
% \let\openright=\cleardoublepage

%% Použité kódování znaků: obvykle latin2, cp1250 nebo utf8:
\usepackage[utf8]{inputenc}

%% Ostatní balíčky
\usepackage{graphicx}
\usepackage{amsthm}
\usepackage{pgfplots}
\usepackage{float}

%% PGF Number format
\pgfkeys{
  /pgf/number format/precision=3,
  /pgf/number format/fixed,
  /pgf/number format/fixed zerofill=true
}

%% Balíček hyperref, kterým jdou vyrábět klikací odkazy v PDF,
%% ale hlavně ho používáme k uložení metadat do PDF (včetně obsahu).
%% POZOR, nezapomeňte vyplnit jméno práce a autora.
\usepackage[ps2pdf,unicode]{hyperref}   % Musí být za všemi ostatními balíčky
\hypersetup{pdftitle=Modular Objective-C Run-Time}
\hypersetup{pdfauthor=Bc. Kryštof Váša}

%%% Drobné úpravy stylu

% Tato makra přesvědčují mírně ošklivým trikem LaTeX, aby hlavičky kapitol
% sázel příčetněji a nevynechával nad nimi spoustu místa. Směle ignorujte.
\makeatletter
\def\@makechapterhead#1{
  {\parindent \z@ \raggedright \normalfont
   \Huge\bfseries \thechapter. #1
   \par\nobreak
   \vskip 20\p@
}}
\def\@makeschapterhead#1{
  {\parindent \z@ \raggedright \normalfont
   \Huge\bfseries #1
   \par\nobreak
   \vskip 20\p@
}}
\makeatother

% Toto makro definuje kapitolu, která není očíslovaná, ale je uvedena v obsahu.
\def\chapwithtoc#1{
\chapter*{#1}
\addcontentsline{toc}{chapter}{#1}
}

\begin{document}

% Trochu volnější nastavení dělení slov, než je default.
\lefthyphenmin=2
\righthyphenmin=2

%%% Titulní strana práce

\pagestyle{empty}
\begin{center}

\large

Charles University in Prague

\medskip

Faculty of Mathematics and Physics

\vfill

{\bf\Large MASTER THESIS}

\vfill

\centerline{\mbox{\includegraphics[width=60mm]{img/logo.eps}}}

\vfill
\vspace{5mm}

{\LARGE Bc. Kry\u{s}tof V\'{a}\u{s}a}

\vspace{15mm}

% Název práce přesně podle zadání
{\LARGE\bfseries Modular Objective-C Run-Time}

\vfill

% Název katedry nebo ústavu, kde byla práce oficiálně zadána
% (dle Organizační struktury MFF UK)
Department of Distributed and Dependable Systems

\vfill

\begin{tabular}{rl}

Supervisor of the master thesis: & Martin D\u{e}ck\'y \\
\noalign{\vspace{2mm}}
Study programme: & programme \\
\noalign{\vspace{2mm}}
Specialization: & specialization \\
\end{tabular}

\vfill

% Zde doplňte rok
Prague 2012-13

\end{center}

\newpage

%%% Následuje vevázaný list -- kopie podepsaného "Zadání diplomové práce".
%%% Toto zadání NENÍ součástí elektronické verze práce, nescanovat.

%%% Na tomto místě mohou být napsána případná poděkování (vedoucímu práce,
%%% konzultantovi, tomu, kdo zapůjčil software, literaturu apod.)

\openright

\noindent
Dedication.

\newpage

%%% Strana s čestným prohlášením k diplomové práci

\vglue 0pt plus 1fill

\noindent
I declare that I carried out this master thesis independently, and only with the cited
sources, literature and other professional sources.

\medskip\noindent
I understand that my work relates to the rights and obligations under the Act No.
121/2000 Coll., the Copyright Act, as amended, in particular the fact that the Charles
University in Prague has the right to conclude a license agreement on the use of this
work as a school work pursuant to Section 60 paragraph 1 of the Copyright Act.

\vspace{10mm}

\hbox{\hbox to 0.5\hsize{%
In ........ date ............
\hss}\hbox to 0.5\hsize{%
signature of the author
\hss}}

\vspace{20mm}
\newpage

%%% Povinná informační strana diplomové práce

\vbox to 0.5\vsize{
\setlength\parindent{0mm}
\setlength\parskip{5mm}

Název práce:
Název práce
% přesně dle zadání

Autor:
Jméno a příjmení autora

Katedra:  % Případně Ústav:
Název katedry či ústavu, kde byla práce oficiálně zadána
% dle Organizační struktury MFF UK

Vedoucí diplomové práce:
Jméno a příjmení s tituly, pracoviště
% dle Organizační struktury MFF UK, případně plný název pracoviště mimo MFF UK

Abstrakt:
% abstrakt v rozsahu 80-200 slov; nejedná se však o opis zadání diplomové práce

Klíčová slova:
% 3 až 5 klíčových slov

\vss}\nobreak\vbox to 0.49\vsize{
\setlength\parindent{0mm}
\setlength\parskip{5mm}

Title:
% přesný překlad názvu práce v angličtině

Author:
Jméno a příjmení autora

Department:
Název katedry či ústavu, kde byla práce oficiálně zadána
% dle Organizační struktury MFF UK v angličtině

Supervisor:
Jméno a příjmení s tituly, pracoviště
% dle Organizační struktury MFF UK, případně plný název pracoviště
% mimo MFF UK v angličtině

Abstract:
% abstrakt v rozsahu 80-200 slov v angličtině; nejedná se však o překlad
% zadání diplomové práce

Keywords:
% 3 až 5 klíčových slov v angličtině

\vss}

\newpage

%%% Strana s automaticky generovaným obsahem diplomové práce. U matematických
%%% prací je přípustné, aby seznam tabulek a zkratek, existují-li, byl umístěn
%%% na začátku práce, místo na jejím konci.

\openright
\pagestyle{plain}
\setcounter{page}{1}
\tableofcontents

%%% Jednotlivé kapitoly práce jsou pro přehlednost uloženy v samostatných souborech
\chapter*{Introduction}
\addcontentsline{toc}{chapter}{Introduction}

Every developer probably has a favorite programming language they feel the most comfortable writing code in. Mine has been Objective-C for a long time now, which made me wonder about the technical background of the language.

I asked myself how does it really work - what is the path from a source code file to a running application. In real life all paths are supposed to lead to Rome. In my case, all paths led to something called \emph{Objective-C run-time}.

While there already are existing Objective-C run-times, a closer look at them shows their weaknesses - mostly their dependencies or assumptions about the underlying OS and environment, in which they are supposed to be running. Be it the POSIX layer, minimal object size, or GCC-specific C language extensions, all of these dependencies may present an obstacle when trying to compile such a run-time for a non-typical operating system (for example an experimental one), or when trying to use the Objective-C run-time in the kernel code.

This thesis analyzes source codes of existing versions of Objective-C run-time, their limitations or requirements for compilation. Result of this work is a prototype of a very flexible run-time in terms of modularity - the run-time is very bare feature-wise, yet offers ways to be easily extended. For example, the run-time includes no support for categories, yet a class extension is included which lets you add support for it. Also, the run-time allows easy configurability - for instance, a single-threaded application is able to turn off locking of internal structures without affecting stability, yet performance can be improved (with each message sent\footnote{In Objective-C method calls are called messages being sent to objects, just like in Smalltalk.}, a lock can be potentially locked when the method implementation is not cached). This may save quite a few syscalls.

The resulting run-time has virtually zero dependencies on the underlying operating system as well as the compiler. The whole run-time is written to be C89 compatible and as a proof of its flexibility, the run-time has not only been run on common systems such as OS X, Windows XP and Linux. It has also been successfully run on more 'exotic' systems, such as Windows 3.11 and Kalisto HeSiVa, an experimental operating system written by me and my colleagues Jiří Helmich and Jan Široký for an Operating Systems course at school.

\chapter{What Is a Run-Time?}

With every object-oriented language a question arises - what makes an object object - how is it represented in the memory. In its core it is a structured piece of memory that could be representable by a structure in the C language. In particular, an object in Objective-C is defined as a structure, whose first field is a so-called \verb=isa= pointer, pointing to the object's class.

With this in mind, a new question pops up - what are methods and how are the method calls performed?

In languages, such as C, it is already known which code should be executed when you call a particular function. The linker then links function calls directly to the address the function resides at. 

Imagine an object-oriented language, where each class had a compile-time known number of methods, these methods could be a part of the class structure, so each method call would consist of just reaching for the correct function pointer within the object's class structure.

This gets more complicated once we take class hierarchy into consideration. Each class can override methods of its parent class, which can be solved by using the same function 'slot'.

Let us demonstrate this on an example written in pseudo-code:

\begin{figure}[H]
  \begin{verbatim}
/** Memory representation of an object. */
struct object_structure {
  class_structure_t *class;
  
  /** Other fields follow. */
} object_structure_t;

/** Memory representation of a class. */
struct class_structure {
  struct class *super_class;
  
  /** Function array of unknown size */
  void*(**functions)(object_structure_t *, ...);
} class_structure_t;
  \end{verbatim}
  \centering{}
  \caption{Structures representing an object and a class.}
  \label{fig:imaginary_lang_structs}
\end{figure}


\begin{figure}[H]
  \begin{verbatim}
class Class1 {
  void method1();
  int method2();
}

class Class2 extends Class1 {
  override void method1();
}
  \end{verbatim}
  \centering{}
  \caption{Declaring two classes.}
  \label{fig:imaginary_lang_classes}
\end{figure}

In this example, \verb=method1= of \verb=Class1= would reside in \verb=Class1->functions[0]=, \verb=method2= in \verb=Class2->functions[1]=, \verb=method1= of \verb=Class2= would reside in \newline{}\verb=Class2->functions[0]=. Hence if \verb=my_obj.method1()= were to be called, it could get translated into \verb=my_obj->class->functions[0]()=.

This, however, brings up a few issues:

\begin{itemize}
  \item{\bf{Binary compatibility}} Linking against a framework which has \verb=method1= in \verb=functions[0]= does not mean the next version will do too - it is the compiler that would decide the order of the functions in the \verb=functions= field.
  \item{\bf{No duck-typing}} If \verb=Class2= is not a subclass of \verb=Class1=, yet implements a method with the same signature - name, types - this would not work either.
  \item{\bf{Extending classes at run time}} All methods must be known at compile time. No methods may be added using e.g. categories like in Objective-C.
\end{itemize}

The solution to these issues is to move to a dynamic dispatch - a mechanism that finds the correct method implementation for that particular object. In this example case, instead of the \verb=functions= field on the class structure, there would be a \verb=methods= field. This methods field would contain a list of methods and the dispatch would look at the name of the method to be invoked, look up the correct implementation and call it. It must be obvious that such a mechanism has its performance cost as calling methods on objects is one of the most common tasks in object-oriented programming.

The \emph{run-time}, being responsible for this dispatch among other tasks, must hence perform this look-up as fast as possible. Other responsibilities of the run-time are to provide a reflection API to constructs of that particular language. For example, listing classes and their methods, exchanging method implementations, listing ivars, etc.


\section{Existing Objective-C Run-Time Implementations}

There are three available Objective-C run-time implementations (to my knowledge):

\begin{itemize}
  \item{\bf{Apple Run-Time}} This run-time is provided by \emph{Apple} and is used in its OS X and iOS systems - there are slight differences between the iOS and OS X versions of the run-time (e.g.\ iOS doesn't support garbage collection and only the new 2.0 run-time is available). Within this thesis, when talking about Apple's implementation of the run-time, the OS X version will be the one talked about.
  \item{\bf{GCC Run-Time}} Also called GNU, or GNUstep run-time, this run-time is provided with the GCC compiler. The naming is a little confusing - the run-time is bundled with the GCC compiler, yet is often referred to as GNU run-time. In 2009, a fork of the run-time's repository has been created to add the newest run-time features and remove legacy code. This fork is called GNUstep run-time. Because both run-times are similar in many ways, both run-times are referred to as the GCC run-time within this thesis.
  \item{\bf{\'{E}toil\'{e}}} \'{E}toil\'{e} is an experimental run-time written by David Chisnall as a part of his paper and has been an inspiration to many improvements in the GNUstep run-time\footnote{http://www.jot.fm/issues/issue\_2009\_01/article4/index.html}.
\end{itemize}


\chapter{Apple's Implementation}

A lot of Apple's source code is \href{''http://www.opensource.apple.com''}{open-sourced} and this includes its implementation of the Objective-C run-time among others.

First thing that comes across mind when studying Apple's source codes is its history - there's a lot of historical code here that ensures backward compatibility. You can find, for example, some code from the NeXT era and a port to Windows - NeXT had a set of APIs called OpenStep (which is the predecessor of today's Cocoa framework on OS X), which was written in Objective-C and was aiming to run on virtually any reasonable system\footnote{http://en.wikipedia.org/wiki/OpenStep}. Also, some of Apple's own software for Windows, such as Safari, was written in Objective-C, hence the run-time needed to be compilable under Windows as well\footnote{http://jongampark.wordpress.com/2009/02/24/safari-4-beta-for-windows/}.

While the legacy code is understandable as it is required to maintain binary compatibility of all existing binaries and the portability is generally an objective of this work, neither is done in a very clean fashion.

All code is duplicated, one part serving compilation of the old run-time and the second part serving compilation of the new run-time. For example, you can find files such as \verb=objc-runtime-old.m= and \verb=objc-runtime-new.mm= (note the \verb=.mm= file extension on the new run-time file suggesting presence of C++, which will be described later).

\section{Portability}
The portability is ensured by rather a large \verb=\#ifdef= - \verb=\#else= - \verb=#endif= statements that use macros and static inline functions to define aliases to some OS X-specific functions on Windows, such as \verb=malloc_zone_malloc= or even POSIX-specific functions, like \verb=issetguid=. I will elaborate on these two examples.

\subsection{Example 1 - malloc}
On OS X, the malloc function has an extension to support memory zones\footnote{http://developer.apple.com/library/Mac/\#documentation/Darwin/Reference/

ManPages/man3/malloc\_zone\_malloc.3.html} - this way you can create multiple zones (which are really just heaps), which can get destroyed entirely at once. It has a very limited usage nowadays, but back in the day when computers had only a little memory, it was useful to allocate temporary objects (e.g.\ during a specific calculation) in its own zone, freeing it as a whole when you were done with these objects. For example, the NSMenu, which is a class representing a menu on OS X, has a class method +menuZone which returns a zone that is used for menu allocations - considering that a menu is an UI element that gets displayed on the screen for only a short period of time, it is a good idea to store all memory used to represent it in a separate zone, freeing all the memory when the menu is dismissed and hence preventing memory fragmentation.

As Windows supports only the regular \verb=malloc= function, this had to be solved - and it has been solved rather radically, by defining the \verb=malloc_*= methods as static inline methods that call the Windows API functions\footnote{objc-os.h}:

\begin{verbatim}
static inline void *malloc_zone_malloc(malloc_zone_t z, 
                                              size_t size) { 
    return malloc(size); 
}
\end{verbatim}

While it's functional, if Microsoft ever decided to implement zones into its version of malloc, it would break the code as it would result in duplicate symbol definition. Moreover, there is another function defined in the source code\footnote{objc-class.m}:

\begin{verbatim}
void *_malloc_internal(size_t size) {
    return malloc_zone_malloc(_objc_internal_zone(), size);
}
\end{verbatim}

It would, of course, make much more sense to simply declare these internal functions that would pass deal with the portability issue themselves (i.e.\ moving the \verb=#if=-\verb=#else= construct here, inside each function would result in a more readable code).

\subsection{Example 2 - issetguid}
The other example is how is the \verb=issetguid= function transferred to the Windows environment:

\begin{verbatim}
#define issetugid() 0
\end{verbatim}

\section{Limitations}
There are several limitations that can be found in the code:

\subsection{16B Object Minimum Size}
Apple supplies several large libraries (or as they call it frameworks), that a vast majority of applications build upon. In particular, the CoreFoundation and Foundation frameworks are used basically by every application in the system.

CoreFoundation, although it implements many Objective-C classes as can be easily verified using the class-dump tool\footnote{http://www.codethecode.com/projects/class-dump/}, they are used privately and the framework provides only C exports and headers.

On the other hand, the Foundation framework, its counterpart, is mainly an Objective-C framework, providing object-oriented access to the mainly C-based CoreFoundation.

For example, \verb=CFStringRef= is a pointer to a structure used by the CoreFoundation framework to represent a string. The Foundation framework has a \verb=NSString= class, which does generally the same. In order to prevent unnecessary code duplication as well as data conversions when using CoreFoundation functions (which accept \verb=CFStringRef=) from Objective-C code, toll-free bridging has been introduced.

There is an intricate mechanism behind it, which requires direct support from the Foundation framework (so you cannot create toll-free bridge of other yet-unsupported Foundation classes)\footnote{You can read more about it here: http://www.mikeash.com/pyblog/friday-qa-2010-01-22-toll-free-bridging-internals.html}, but instances of classes that do support toll-free bridging, can be simply casted to their CoreFoundation counterpart and vice versa. Using the \verb=CFStringRef=/\verb=NSString= couple, this code if fully valid:

\begin{verbatim}
NSString *myString = @"Hello World!";
CFStringRef duplicatedString = CFStringCreateCopy(NULL, 
                                    (CFStringRef)myString);
NSString *duplicatedString2 = (NSString*)duplicatedString;
\end{verbatim}

While this is very convenient to every OS X (or iOS) developer, it poses an unexpected limitation: all object instances need to have the same minimum size as CoreFoundation ``objects" - 16 bytes, which is noted in a comment in \verb=objc-class.m= file within the \verb=_class_createInstancesFromZone= function.

\subsection{Dynamic Loader Support}
As it's faster to simply copy the data from the binary than to construct the classes using the run-time functions (on OS X from the \verb=__OBJC= section in particular), Apple's implementation contains a set of functions that are called by the dynamic loader (\verb=dyld=) to load classes from the binary image, link them to their superclasses and to register them with the run-time.

This, however, adds \verb=dyld= to the list of library dependencies, which is transitively required by \verb=libSystem=\footnote{Basic system library on OS X that every application needs to be linked to.} anyway, but it adds dependencies to the code itself.

Also, loading a binary image of the classes introduces an issue with binary compatibility - any change in the internal representation of a class will cause the binary not to launch.

Creating classes manually using the run-time methods ensures binary compatibility, while slows down binary loading significantly and poses a question where should be the class information stored. One option would be to generate a function with a specific name that's called by the dynamic loader, or calling it manually before anything else in the \verb=main= of your program (which might interfere with the \verb=__attribute__((constructor))= functions, though).

\subsection{C++ Influences}
The newest parts of the Objective-C run-time use many C++ features, such as methods on structures, some C++ classes, e.g. \verb=vector= and the run-time tries to unify Objective-C and C++ exceptions using the unwind library.

While this may help to clean up the code a little, it adds additional dependencies on C++ libraries.

\section{Summary}
Apple's implementation is riddled with a lot of obscure code and other very specific details - for example, when the class images are stored in the \verb=__OBJC= binary section, the superclass field is pointing to a string containing the name of that superclass - when the run-time is connecting the images, it's casting the superclasses pointer to char*, to read the superclass name, which is not a very clean solution (for example, the \'Etoil'e run-time uses two differently named structures for a cleaner cast). Moreover, the documentation is very brief, not every function has a description of what it exactly does, or only has a very short note that it is used from a different function somewhere else in the code. As will be described below, the GCC version of the run-time is much better documented in this matter.

\chapter{Apple's Implementation}

A lot of Apple's source code is \href{''http://www.opensource.apple.com''}{open-sourced} which includes their implementation of the Objective-C run-time among others.

First thing that comes across mind when studying Apple's source codes is its history - there's a lot of historical code that ensures backward compatibility. You can find, for example, some code from the NeXT era and a port to Windows - NeXT had a set of APIs called OpenStep (which is the predecessor of today's Cocoa framework on OS X), which was written in Objective-C and was aiming to run on virtually any reasonable system\footnote{http://en.wikipedia.org/wiki/OpenStep}. Also, some of Apple's own software for Windows, such as Safari, is written in Objective-C, hence the run-time needed to be compatible with Windows as well\footnote{http://jongampark.wordpress.com/2009/02/24/safari-4-beta-for-windows/}.

While the legacy code is understandable as it is required to maintain binary compatibility of all existing binaries and the portability is generally an objective of this work, neither is done in a very clean fashion.

The code that was used by the run-time before introduction of Objective-C 2.0 in 2007, is referred to as the \emph{old run-time}. The code that is part of the Objective-C 2.0 (and later) if referred to as the \emph{new run-time}.

All code is duplicated, one part serving compilation of the old run-time and the second part serving compilation of the new run-time. For example, you can find files such as \verb=objc-runtime-old.m= and \verb=objc-runtime-new.mm= (note the \verb=.mm= file extension on the new run-time file suggesting presence of C++, which will be described later).

\section{Portability}

The portability of the run-time can be divided in two parts - portability in terms of operating systems and in terms of the CPU architecture.

\subsection{CPU Portability}

Historically, Objective-C has been run on quite a few architectures - Intel x86, and Sun microchips from the NeXT era, PPC, Intel x86-64 and more recently ARM chips.

The \verb=objc_msgSend= function and its relatives are written in assembly and hence has to be rewritten for each new architecture.

The OS portability is ensured by rather a large \verb=#ifdef= - \verb=#else= - \verb=#endif= statement that uses macros and static inline functions to define aliases to some OS X-specific functions on Windows, such as \verb=malloc_zone_malloc= or even POSIX-specific functions, like \verb=issetguid=. Two examples follow.

\subsection{Example 1 - malloc}
On OS X, the malloc function is extended to support memory zones\footnote{http://developer.apple.com/library/Mac/\#documentation/Darwin/Reference/

ManPages/man3/malloc\_zone\_malloc.3.html}. By creating a zone, a new heap is created. The heap can be destroyed entirely at once. It has a very limited usage nowadays, but back in the day when computers had only a little memory, it was useful to allocate temporary objects (e.g.\ during a specific calculation) in its own zone, freeing it as a whole when you were done with these objects. For example, the \verb=NSMenu=, which is a class representing a menu on OS X, implements a class method \verb=+menuZone= which returns a zone that is used for menu allocations - considering that a menu is an UI element that gets displayed on the screen for only a short period of time, it is a good idea to store all memory used to represent it in a separate zone, freeing all the memory when the menu is dismissed and hence preventing memory fragmentation. Unfortunately, when releasing a chunk of memory like this at once, the \verb=-dealloc= method doesn't get called, which could result in memory leaks. The zones were also supposed to be used to add garbage collection support to Objective-C back in the NeXT days, which never happened, though.

As Windows supports only the regular \verb=malloc= function, this had to be solved - and it has been solved rather radically, by defining the \verb=malloc_*= functions as static inline functions that call the Windows API functions. The following lines of code may be found in file \verb=objc-os.h=:

\begin{verbatim}
static inline void *malloc_zone_malloc(malloc_zone_t z, 
                                              size_t size) { 
    return malloc(size); 
}
\end{verbatim}

While it is functional, if Microsoft ever decided to implement zones into its version of malloc, it would break the code as it would result in duplicate symbol definition. Moreover, there is another function defined in the source code file \verb=objc-class.m=:

\begin{verbatim}
void *_malloc_internal(size_t size) {
    return malloc_zone_malloc(_objc_internal_zone(), size);
}
\end{verbatim}

It would make much more sense to simply declare functions that would deal with the portability issue themselves, i.e.\ moving the \verb=#if=-\verb=#else= construct inside each function which would result in a more readable code.

\subsection{Example 2 - issetguid}
The other example is how is the \verb=issetguid= function ported to the Windows environment:

\begin{verbatim}
#define issetugid() 0
\end{verbatim}

\section{Limitations}
There are several limitations that can be found in the code that may present an obstacle when compiling the run-time on a new platform.

\subsection{16B Object Minimum Size}
Apple supplies several large libraries (or as they call them, frameworks), that a vast majority of applications build upon. In particular, the CoreFoundation and Foundation frameworks are used basically by every application in the system (if not directly, then at least indirectly).

CoreFoundation, although it implements many Objective-C classes as can be easily verified using the class-dump tool\footnote{http://www.codethecode.com/projects/class-dump/}, only provides C exports and headers, with the Objective-C classes being used privately.

On the other hand, the Foundation framework, CoreFoundation's counterpart, is mainly an Objective-C framework, providing object-oriented access to the mainly C-based CoreFoundation and officially allows access to some Objective-C classes declared in CoreFoundation.

For example, \verb=CFStringRef= is a pointer to a structure used by the CoreFoundation framework to represent a string. The Foundation framework has a \verb=NSString= class, which does generally the same. In order to prevent unnecessary code duplication as well as data conversions when using CoreFoundation functions (which accept \verb=CFStringRef=) from Objective-C code, toll-free bridging has been introduced.

There is an intricate mechanism behind it, which requires direct support from both the Foundation and CoreFoundation frameworks (so one cannot create toll-free bridge of other (Core)Foundation classes that aren't bridged yet himself or herself)\footnote{You can read more about it here: http://www.mikeash.com/pyblog/friday-qa-2010-01-22-toll-free-bridging-internals.html}, but instances of classes that do support toll-free bridging, can be simply casted to their CoreFoundation counterparts and vice versa. Using the \verb=CFStringRef=/\verb=NSString= couple, this code if fully valid:

\begin{figure}[H]
\begin{verbatim}
NSString *myString = @"Hello World!";
CFStringRef duplicatedString = CFStringCreateCopy(NULL, 
                                    (CFStringRef)myString);
NSString *duplicatedString2 = (NSString*)duplicatedString;
\end{verbatim}
  \centering{}
  \caption{An example of toll-free bridging.}
  \label{fig:toll_free_bridging}
\end{figure}

While this is very convenient to every OS X (or iOS) developer, it poses an unexpected limitation: all object instances need to have the same minimum size as CoreFoundation ``objects" - 16 bytes, which is noted in a comment in \verb=objc-class.m= file within the \verb=_class_createInstancesFromZone= function. True that 16 bytes isn't that much in these days, but given that you could fit 4 regular objects into one on a 32-bit computer, the space adds up.

\subsection{Dynamic Loader Support}
As it is faster to simply copy the class data from the binary than to construct the classes using run-time functions (on OS X from the \verb=__OBJC= section in particular), Apple's implementation contains a set of functions that are called by the dynamic loader (\verb=dyld=) to load classes from the binary image, link them to their superclasses and to register them with the run-time.

This, however, adds \verb=dyld= to the list of library dependencies, which is transitively required by \verb=libSystem=\footnote{Basic system library on OS X that every application needs to be linked to.} anyway, but it adds dependencies to the code itself.

More portable code would declare structures to be passed to the run-time when loading a binary, and the dynamic loader would pass those to the run-time. Instead, the dynamic loader passes the binary image header pointer to the run-time, making the run-time crawl through the binary image headers for Objective-C data.

In general, the dynamic loader support in Apple's run-time poses a question, which party is responsible exactly for what part of the binary loading. It also puts a light on the binary compatibility question. If the internal structures representing a class change, the binary will not launch.

Creating classes manually using the run-time methods ensures binary compatibility, while slows down binary loading significantly and makes one ask where should be the class information stored. One option would be for the compiler to generate a function with a specific name that would get called by the dynamic loader, or called it manually before anything else in the \verb=main= of your program (which might interfere with the \verb=__attribute__((constructor))= functions, though).

\subsection{C++ Influences}
The newest parts of the Objective-C run-time use many C++ features, such as methods on structures, some C++ classes, e.g. \verb=vector= and the run-time tries to unify Objective-C and C++ exceptions using the \verb=unwind= library.

While this may help to clean up the code a little, it adds additional dependencies on C++ libraries.

\section{Summary}
Apple's implementation is riddled with a lot of obscure code and other very specific details - for example, when the class images are stored in the \verb=__OBJC= binary section, the superclass field is pointing to a string containing the name of that superclass - when the run-time is connecting the images, it is casting the superclasses pointer to char*, to read the superclass name, which is not a very clean solution. For example, the \'Etoil'e run-time uses two differently named structures for a cleaner cast. 

Moreover, the documentation is very brief, not every function has a description of what it exactly does, or only has a very short note that it is used from a different function somewhere else in the code. As will be described below, the GCC version of the run-time is much better documented in this matter.

\chapter{\'Etoil\'e}
\section{What is \'Etoil\'e}

\'Etoil\'e is an experimental run-time written by David Chisnall, a research assistant in computer science at Swansea University. It has been written as a part of his paper called \emph{A Modern Objective-C Runtime}\footnote{http://www.jot.fm/issues/issue\_2009\_01/article4/index.html}. Even though it is not a real-world run-time (i.e.\ no compiler supports it in its full extent and has been claimed to be mostly experimental), it has introduced quite a few interesting ideas how to speed up the method lookup time as well as make the run-time more generic.

Unlike the previously described run-times, the \'Etoil\'e run-time tackles the task of providing a run-time from a totally different angle. While the other implementations simply aim to create a traditional Objective-C run-time, where Apple chooses to keep almost all of the original API for its Objective-C 2.0, \'Etoil\'e run-time tries to create a very generic run-time that could be used with many other languages as well, which would result in a very easy language bridging - a mechanism, where objects from one language can be interacted with from another language and vice versa (even though deprecated, Apple was supplying a Java - Objective-C bridge for quite some time\footnote{http://cocoadevcentral.com/articles/000024.php}).

This task, however, required to start from scratch and leave all compatibility behind. As the author notes, the run-time API itself has never been standardized, unlike the language, so a person shouldn't rely that much on it. Hence all the functions, such as \verb=objc_msgSend=, \verb=objc_getClass=, or \verb=sel_registerName= are not available in this run-time. The data structures are modified, or completely missing - for example \verb=SEL= is defined as a \verb=uint32_t=, which is a hash for an internal representation of the selector, which is a \verb=objc_selector= structure, which contains name and a type string, like in GCC's run-time. 

The source code of this run-time is much shorter than the other implementations - according to the author the run-time is just 15\% of the size of GCC's implementation, however, this run-time doesn't currently provide some of the modern features of run-times. The source code includes many macros that are really to be generated by the compiler as well - the speed of the run-time relies on the compiler a lot as well.

\section{Slots}
Instead of defining methods, the more generic approach is to define slots. A slot is the basic type for message lookup - a structure containing five fields: \verb=int offset=, \verb=IMP method=, \verb=char *types=, \verb=void *context= and \verb=uint32_t version=. This allows the run-time to store both properties and methods using a single structure: as the properties, introduced in Objective-C 2.0 run-time, are just wrappers around automatically generated getter and setter methods, this approach allows to define a property simply using a slot that has a defined \verb=offset= field, which serves as a number of bytes after the beginning of the object is that variable stored. By setting the \verb=method= field of the slot, the slot works like a regular method.

As I have mentioned, the run-time was built to meet needs of other languages as well, for example JavaScript (and other prototype-based languages), where you can dynamically add variables to an object and inherit not just from a class, but from another object as well. This can be done by adding a slot, which will then hold the value as well - in the \verb=context= field of the structure.

Whenever the structure gets updated (\verb=IMP= is changed, etc.), the version is increased. This is important later on for caching.


\section{Inline Caching}
Using monomorphic or polymorphic inline caching, the author of \'Etoil\'e run-time was able to achieve impressive speeds, reducing the method call time to only twice the time of a pure C function call, even faster than the C++ method calls are.

With every dynamic object-oriented language, a question arises, how to fetch the function that actually implements a method. This lookup function is usually the critical part of the run-time's performance. As the lookup can be expensive (climbing the whole class hierarchy), all of the run-times described here use some caching mechanisms.

Imagine your class \verb=FCButton=, which is a subclass of \verb=NSButton=, which is a subclass of \verb=NSControl=, \verb=NSView=, \verb=NSResponder=, \verb=NSObject=. If the dispatch had to lookup each method, for example the very commonly used \verb=retain= and \verb=release=, which are implemented most likely just in the \verb=NSObject= root class, the lookup function would have to climb the whole class hierarchy, until it found the class that implements the method.

For this reason, caching has been introduced. Both Apple's and GCC's run-times have a cache - when the user wants to call a method, the run-time needs to look up the function pointer. First, it looks into the cache (which is usually fast) and returns the \verb=IMP=, if it has been found, otherwise looks up the function pointer in the class hierarchy and saves it into the cache for further use.

An issue here is that when the implementation of a method in some class is changed, all cache entries with the original function pointer need to be removed. The same applies when a new module has been loaded with a class category - a category may replace an already-existing method of that class.

But even this lookup costs something - there's at least one C function call to fetch the data from the cache, plus the actual fetching from the cache. The approach of \'Etoil\'e run-time is for the compiler to generate inline caches.

Every time a method is supposed to be called, the following lookup macro is applied. This example uses the monomorphic cache. Unless you assume the same message will sent to objects of different classes, it is well sufficient. The polymorphic cache may be useful for class clusters (for example, NSString - which can be of multiple different classes, depending whether it's a constant string, created by run-time, etc.), or if you are expecting subclasses of the variable type class to be passed as well.

\begin{verbatim}
#define SLOT_LOOKUP_MIC(obj, sel_name, sel_types, sender, action)\
do\
{\
    static SEL selector = 0;\
    struct objc_slot * slot;\
    if(selector == 0)\
    {\
        selector = lookup_typed_selector(sel_name, sel_types);\
    }\
    static __thread struct inline_cache_line cache;\
    if(cache.slot != NULL \
        &&\
        cache.type == (id)obj->isa\
        &&\
        cache.version == cache.slot->version)\
    {\
        slot = cache.slot;\
    }\
    else\
    {\
        id object = (id)obj;\
        slot = slot_lookup(&obj, selector, sender);\
        if(obj == object)\
        {\
            cache.version = slot->version;\
            cache.slot = slot;\
            cache.type = obj->isa;\
        }\
    }\
    struct objc_call call = { slot, selector, sender };\
    action\
} while(0)
\end{verbatim}


Inline caches are static \verb=__thread= structures (which as has been noted above, is a GCC extension, that needs support from the OS), so that each thread has its own cache. This static cache is created for every place in the code, where a method is called. The selector itself is cached, so no lookup is needed for it either. Then the cache is checked if it is filled with this object's class (i.e. this place in the code has been visited before, but the object has been of another class before - a polymorphic cache should have been used here instead) and the version matches the slot's version - if the slot has been modified since, the cache is invalidated and filled freshly.

\section{Message Sending}

The traditional message sending, where the method gets translated into a function with the first parameter \verb=self= and the second one \verb=_cmd= has been abandoned for a slightly more complicated, yet more flexible call:

\begin{verbatim}
typedef struct objc_call {
  SLOT slot;
  SEL selector;
  id sender;
} *CALL;
#define _cmd (_call->selector)
id method(id self, CALL _call, ...);
\end{verbatim}

This extends the original simple \verb=_cmd= of type \verb=SEL=: it adds more context to the call, which can be used to implement a private inner class, for example, which determines with each call that the \verb=sender= is indeed its parent class.

\section{Tags}

Each object has a couple of tags that the run-time installs by default - additional fields that aren't declared by the object itself. In particular, each object has its own mutex, refcount (number of references to the object - used for reference-based memory management), size, slots and lookup function.

As each object has its own mutex, unlike Apple's and GCC's run-times, where a pool of locks is used to minimize memory usage as well as the cost related to initializing a recursive mutex, even though some systems have lazy-initialization available. In Apple's and GCC's run-time, a hash is computed for each object, which determines which lock should be used. If it has already been taken by another object, another lock is used. The pool is claimed to be big enough, though, that in regular situations no contention is encountered. \'Etoil\'e's approach, however, decreases the contention to none at all.

Each object also has its own dispatch table, unlike the other run-times. This allows to change the dispatch object-based, i.e. swapping method implementations for just one object, adding a method for just one object, etc.

This leads to the last interesting tag, which is the lookup function. A lookup function is a function that goes through the dispatch table of the object and looks for the implementation of a method you want to call. Here, you can change the lookup function simply on an per-object basis, which allows to implement proxy objects very efficiently, making proxy calls virtually the same speed as direct calls (by setting the lookup function to a function that looks directly through the target class' methods).

\section{Metaclasses}

Another difference between the \'Etoil\'e run-time and the others is the lack of meta classes. The class methods are simply put into a dispatch table (in this case only a re-defined sparse array) directly on the class. In the traditional run-time, the instance methods have a dispatch table on the class, class methods are instance methods of its meta class. Here, the dispatch table of instance methods is installed directly onto the object, while the class methods have a dispatch table on the class itself. Hence no meta classes are needed in this run-time.


\chapter{Modular Run-Time}
\section{Objectives}

The main objective of this work is to create a prototype of a run-time that will have no dependencies on any external libraries, not even the POSIX ones, allowing the developer to set up his own run-time from different modules.

Also, no constructs that require any OS support (dynamic loader, standard OS libraries, etc.) will be used. This means, that constructs such as \verb=static\ __thread= (which creates a new copy of the variable declared for each thread created) or \newline{} \verb=__attribute__((constructor))= (which marks the function to be called by the dynamic loader right after the binary has been loaded) can be used either.

This will, however, requires the user to call some run-time initializing functions manually at the beginning of the program's execution, as they won't get called by the dynamic loader. On systems that do support the constructor functions, this can be easily avoided by compiling the run-time with an additional file which will declare and implement a constructor, as can be seen in the \verb=posix.c= file, which is included with the run-time as an extra, which demonstrates how to add support for your operating system.

In a real-world scenario, of course, we cannot require that each user supplies all these functions first thing in the \verb=main= function. For example, on \verb=*nix=\footnote{Most Unix systems, cannot be guaranteed that this will work on all of the systems.} systems, this can be easily solved by compiling the library with a single \verb=.c= file which has a function populating the setup structure with system-specific functions. This function can then either be called by the dynamic loader (which requires support from the dynamic loader), or can have a GCC-specific \verb=__attribute__((constructor))= if the system supports it. If neither is an option, then the user really needs to call this system-specific function to populate the run-time setup and then initialize the run-time, if the system-specific function hasn't done so. An example of this can be found in the \verb=posix.c= file which is bundled with the run-time as an example.


\subsection{Example 1: Lock-less run-time}
In a single-threaded application (or applications, where you know that Objective-C code will be used only in one thread), there is no need for any locking whatsoever. All of the existing implementations require some locking, even though they are trying to use lock-free structures, such as sparse arrays that do not support deleting.

But even so, any \verb=@synchronized(obj)= code is translated to actually lock a mutex associated with \verb=obj=, be it either a mutex from a lock pool in the traditional run-times, or a mutex that is associated just with \verb=obj= in the \'Etoil\'e run-time. This can speed up both loading of the application and code execution.

\subsection{Example 2: Kernel usage}

To get the existing run-times working in a kernel of an operating system might be tricky, depending on how much is the kernel POSIX-compatible. But even so, the \verb=malloc= functions and others usually are just wrappers around kernel allocators, which slows down allocation of all structures within the run-time.

Using the modular run-time it will be possible to change the allocator with a simple function-pointer assignment.


\subsection{Example 3: Benchmarking}

The modularity that will be introduced by this work will allow anyone to explore changes in the speed of the run-time simply by changing internal data structures used to hold the class list, selector list and caching. This may help the future development of the run-time.


\section{Run-time Setup}

At the core of the run-time, there is a structure \verb=objc_runtime_setup_struct objc_setup= - while the user has no direct access to this structure, the run-time itself can access it directly for the sake of the speed, to eliminate unnecessary function calls that would serve just as proxy calls.\footnote{The user doesn't have a direct access to the structure only as a security precaution so that the structure cannot be modified from the outside during the program execution. Of course, the user can easily get the setup using the 'objc\_runtime\_get\_setup' function, which returns the current setup structure. The user can then cache this structure for performance.}

You can view the structure here:

\begin{verbatim}
  // TODO: paste the final structure when available
typedef struct {
  objc_abort abort;
  objc_allocator allocator;
  struct {
    objc_class_holder_creator creator;
    objc_class_holder_destroyer destroyer;
    objc_class_holder_lookup lookup;
  } class_holder;
} objc_runtime_setup_struct;
\end{verbatim}

This allows the modularity of the run-time. One can, at the beginning of his program, modify all of those pointers using setter functions declared in \verb=runtime.h=. After the runtime has been initialized, however, the whole structure is sealed from changes using those setter functions to prevent any data corruption (as some data structures may be already initialized, changing these functions could lead to bad memory access, etc.). Using any of the setter functions after the run-time has been initialized will lead to a crash of the program (using the abort function supplied).

A description of each of those functions follows:

\subsection{Abort}

It is commonly said that there is no software without bugs - hence it is sometimes necessary to abort the program execution if the program gets into an inconsistent point, or a point where illegal arguments are passed to the environment (for example, setting the run-time setup structure to NULL).

\subsection{Allocator, Deallocator, Reallocator and Zero-allocator}

As the run-time needs new memory for dynamic object creation, an allocator is needed. All existing run-times use \verb=malloc=, which, however, ties them to Unix-based operating systems.

The Modular Run-Time lets you specify your own allocator, which can be, for example, just a wrapper around \verb=malloc= adding some debug logging, or a completely different allocator, e.g.\ a kernel allocator as has been mentioned before. It can also be just the \verb=malloc= function itself, as the function type takes just one argument - the size of the memory required and returns a \verb=void*= pointer.

Sometimes, the memory acquired should be filled with zeros - just as the \verb=calloc= function would do, however, it is called \verb=objc_zero_allocator\ zero_allocator= in this run-time. The \verb=calloc= function can be assigned, however.

For dynamically growing structures, reallocation may be needed - hence the inclusion of the reallocating function. It is pointer-compatible with the \verb=realloc= function.

When the run-time is done with memory it has allocated, it will deallocate it using the deallocate function, which has the same signature as the POSIX \verb=free= function.

\subsection{Arrays}

A lot of the code of the traditional run-times is riddled with code that is takes care of consistency of dynamically growing arrays (or rather arrays of arrays) - there's a lot of duplicate code for functions related to lists of methods, protocols, etc.\ on each class.

This run-time declares a \verb=objc_array= type, which is just a retyped \verb=void*=, but can be implemented in any possible way. The run-time includes a working implementation of such an array and installs these function pointers at initialization, unless other pointers are provided.

The default implementation is a structure that contains two \verb=unsigned int=s - one that remembers the allocated capacity of the array, the second one pointing to an index in the actual allocated array holding the items. If the structure is initialized as lock-able, then a mutex is allocated as well.


\subsection{Class Holder}

The run-time needs to keep track of all classes that are registered with it. To do so, it needs to keep a list in some data structure. As the run-time is designed to be flexible, it doesn't matter what data structure at all.

There is a defined data type \verb=objc_class_holder= which, however, is again only a retyped \verb=void*=. Using the run-times function pointers \verb=class_holder.creator=, \verb=class_holder.destroyer=, \verb=class_holder.lookup= and TODO, the run-time creates the data structure required, has means to destroy it, lookup a class according to its name and TODO.

The run-time provides a default implementation of this structure, a simple hash table which resembles slightly \verb=NXMapTable= from Apple's run-time, but removes all dependencies and is made slightly less complicated.

\section{Class}

At the core of the run-time lies a structure representing a class. Like the \'Etoil\'e run-time, this run-time doesn't provide a class pair - a class and its meta-class, but only a single class object that contains class methods as well. While it abandons the Smalltalk similarities, it provides greater flexibility, allowing the Objective-C class structure to be used for other languages as well.

\subsection{Structure of a Class}

The class structure begins in the traditional fashion with a pointer to its superclass - the \verb=isa= pointer. This pointer is followed by the class name, class methods and instance methods. Both class methods and instance methods use the \verb=objc_array= structure, which is a list of lists of \verb=Method=s, just like in the other run-times 

// TODO: cache, ivars

\section{Extensibility of Classes}

While the other run-times allow the user to specify extra space when allocating both a class and an object, it is very limiting as if you want to add this extra space, you need to replace the \verb=+alloc= method of \verb=NSObject=, but even so, it will not affect all classes in the run-time as not all objects are subclasses of \verb=NSObject= (e.g.\ \verb=NSProxy=).

This is why the modular run-time has introduced class extensions which allow developers to extend the run-time capabilities dynamically. Class extensions can be compared to the run-time delegates. The class extensions must get registered using the \verb=objc_class_add_extension()= function before the run-time gets initialized.

This function has a single argument: a pointer to a \verb=objc_class_extension= structure:

\begin{verbatim}
  typedef struct _objc_class_extension {
	struct _objc_class_extension *next_extension;


	void(*class_initializer)(Class, void*);
	void(*object_initializer)(id, void*);
	void(*object_deallocator)(id, void*);
	
	unsigned int extra_class_space;
	unsigned int extra_object_space;
} objc_class_extension;
\end{verbatim}

The first field, \verb=next_extension= is a pointer to the next extension as the run-time keeps class extensions in a linked list. The run-time populates this field automatically.

This field is followed by three function pointers:

\begin{itemize}
  \item \verb=class_initialize= - this function gets called each time a class is registered with the run-time using \verb=objc_finish_class()=. This is the place where the extra space, to which a pointer is passed as the second argument, should get populated and initialized.
  \item \verb=object_initializer= - any time an object is allocated, this function gets called in order to initialize the extra space in the object (again, the second argument is a pointer to the extra space within the structure).
  \item \verb=object_deallocator= - before the object gets deallocated, each class extension is provided an option to deallocated any dynamically allocated structures.
\end{itemize}

While it allows to extend the classes with new functionality, for example properties, or ARC (the \verb=object_deallocator= can do all the automatic deallocations), it poses an issue at the compile time - how much extra space should be left? The compiler needs to know the size of the class structure to generate.

One option is to re-allocate all classes with the extra space and copy over all pointers of the class internal structures (it is enough to copy over pointers as the memory will be kept alive). Assuming 1000 classes in a larger application, each class having 64 bytes + those extra bytes, this gives 64 kB of extra memory per application. While this isn't that much nowadays, it would slow down the application launch.

Other option is to dynamically allocate the extra space for each class, so the class structure stays of the same size, with the extra space being outside of the structure itself.


There is an example bundled with the run-time: properties.

\subsection{Properties}

The run-time itself doesn't include support for properties, however, the run-time provides support for extending both the class and the object itself.

// TODO describe


\subsection{Performance}
\chapter{Known Issues}

// if properties were implemented using CE, when and how to install them on the class?

+load
+initialize

\_\_MRConstString - class extensions


\chapter{Future Work}

While the run-time provides already a lot of capability, there are some features of the other run-times that is misses, or could be worked on in the future. This chapter will list a few of these issues and will elaborate how easily this functionality could be added to the modular run-time.

\begin{itemize}
  \item{\bf{Compiling information for class extensions}} As the categories are unknown at the compile time and the space they occupy on the \verb=Class= object is dynamically allocated, how does one compile, for example class categories, and register them with the run-time? One possibility is to compile the categories into prototypes, which is easy, register the class prototype and then add the category to the class. Other alternative is to use the \verb=extra_space= field in the class prototype structure, making it a linked list of extension data prototypes. Each function would have a string identifier which would serve to find the correct extension data prototype. To implement this, all that is needed is to modify the class prototype structure and slightly extend the class registering function.
  \item{\bf{\tt{\+load} method}} In the traditional run-times, when a method gets registered with the run-time (or more likely loaded from a bundle) and it implements the \verb=+load= method (directly that class, not its superclasses!), it gets called. This can be achieved by adding an extension that calls this method within its \verb=class_initializer= function. This however, would require the either for the \verb=+load= method not to use any extensions (as other extensions may not be done with the \verb=class_initializer=s), or somehow make sure it is called as the last. Other solution is to extend the class extension structure with another function pointer, which would get called after the class has been completely initialized and installed into the run-time.
  \item{\bf{\tt{\+initialize} method}} Just like the \verb=+load= method, there is a \verb=+initialize= method that gets called before the first method gets to be invoked on that particular class. This can be easily done by installing a class extension with a look-up function, which always returns \verb=NULL=, but when it gets called for the first time for some class, it looks up and calls (itself) the \verb=+initialize= method. The extension may request additional space on the \verb=Class= structure of size \verb=sizeof(BOOL)=, which would be set to \verb=YES= once the method would get to be called.
  \item{\bf{Static object instances and class extensions}} There aren't many cases where static object instances are necessary, however, for example string literals are the case. Since the compiler doesn't know the additional space requested by the class extensions during compile time, class extensions would fail on \verb=__MRConstString= instances. This could get solved by adding an \verb=extra_space= ivar to the class and check the \verb=isa= pointer in the \newline{}\verb=objc_object_extensions_beginning_for_extension= function - if it indeed is the \verb=__MRConstString= class (a single pointer comparison), then see the \verb=extra_space= field - if it is \verb=NULL=, allocate it (similar as with the \verb=extra_space= on the \verb=Class= objects), otherwise, return a pointer from within that dynamically allocated space. This would mean lazily allocated extra space and the class extensions issue would be solved. Unfortunately, the issue that the object initializer functions wouldn't get called cannot be easily solved, however, these functions should only be used for debugging, since everything should be lazily initialized.
  \item{\bf{Synchronization}} The \verb=@synchronize= construct requires some object-related locking, which can, again, be implemented using a class extension, which can decide itself whether to use a pool of locks, or install a lock on each object.
\end{itemize}



\chapter{Future Work}

While the run-time provides already a lot of capability, there are some features of the other run-times that is misses, or could be worked on in the future. This chapter will list a few of these issues and will elaborate how easily this functionality could be added to the modular run-time.

\begin{itemize}
  \item{\bf{Compiling information for class extensions}} As the categories are unknown at the compile time and the space they occupy on the \verb=Class= object is dynamically allocated, how does one compile, for example class categories, and register them with the run-time? One possibility is to compile the categories into prototypes, which is easy, register the class prototype and then add the category to the class. Other alternative is to use the \verb=extra_space= field in the class prototype structure, making it a linked list of extension data prototypes. Each function would have a string identifier which would serve to find the correct extension data prototype. To implement this, all that is needed is to modify the class prototype structure and slightly extend the class registering function.
  \item{\bf{\tt{\+load} method}} In the traditional run-times, when a method gets registered with the run-time (or more likely loaded from a bundle) and it implements the \verb=+load= method (directly that class, not its superclasses!), it gets called. This can be achieved by adding an extension that calls this method within its \verb=class_initializer= function. This however, would require the either for the \verb=+load= method not to use any extensions (as other extensions may not be done with the \verb=class_initializer=s), or somehow make sure it is called as the last. Other solution is to extend the class extension structure with another function pointer, which would get called after the class has been completely initialized and installed into the run-time.
  \item{\bf{\tt{\+initialize} method}} Just like the \verb=+load= method, there is a \verb=+initialize= method that gets called before the first method gets to be invoked on that particular class. This can be easily done by installing a class extension with a look-up function, which always returns \verb=NULL=, but when it gets called for the first time for some class, it looks up and calls (itself) the \verb=+initialize= method. The extension may request additional space on the \verb=Class= structure of size \verb=sizeof(BOOL)=, which would be set to \verb=YES= once the method would get to be called.
  \item{\bf{Static object instances and class extensions}} There aren't many cases where static object instances are necessary, however, for example string literals are the case. Since the compiler doesn't know the additional space requested by the class extensions during compile time, class extensions would fail on \verb=__MRConstString= instances. This could get solved by adding an \verb=extra_space= ivar to the class and check the \verb=isa= pointer in the \newline{}\verb=objc_object_extensions_beginning_for_extension= function - if it indeed is the \verb=__MRConstString= class (a single pointer comparison), then see the \verb=extra_space= field - if it is \verb=NULL=, allocate it (similar as with the \verb=extra_space= on the \verb=Class= objects), otherwise, return a pointer from within that dynamically allocated space. This would mean lazily allocated extra space and the class extensions issue would be solved. Unfortunately, the issue that the object initializer functions wouldn't get called cannot be easily solved, however, these functions should only be used for debugging, since everything should be lazily initialized.
  \item{\bf{Synchronization}} The \verb=@synchronize= construct requires some object-related locking, which can, again, be implemented using a class extension, which can decide itself whether to use a pool of locks, or install a lock on each object.
\end{itemize}




% Ukázka použití některých konstrukcí LateXu (odkomentujte, chcete-li)
% %%% Ukázka použití některých konstrukcí LaTeXu

\subsection{Ukázka \LaTeX{}u}
\label{ssec:ukazka}

This short subsection serves as an~example of basic \LaTeX constructs,
which can be useful for writing a~thesis.

Let us start with lists:

\begin{itemize}
\item The logo of Matfyz is displayed in figure~\ref{fig:mff}.
\item This is subsection~\ref{ssec:ukazka}.
\item Citing literature~\cite{lamport94}.
\end{itemize}

Different kinds of dashes:
red-black (short),
pages 16--22 (middle),
$45-44$ (minus),
and this is --- as you could have expected --- a~sentence-level dash,
which is the longest.
(Note that we have follwed \verb|a| by a~tilde instead of a~space
to avoid line breaks at that place.)

\newtheorem{theorem}{Theorem}
\newtheorem*{define}{Definition}	% Definice nečíslujeme, proto "*"

\begin{define}
A~{\sl Tree} is a connected graph with no cycles.
\end{define}

\begin{theorem}
This theorem is false.
\end{theorem}

\begin{proof}
False theorems do not have proofs.
\end{proof}

\begin{figure}
	\centering
	\includegraphics[width=30mm]{../img/logo.eps}
	\caption{Logo of MFF UK}
	\label{fig:mff}
\end{figure}


\chapter*{Conclusion}
\addcontentsline{toc}{chapter}{Conclusion}


%%% Seznam použité literatury
%%% Seznam použité literatury je zpracován podle platných standardů. Povinnou citační
%%% normou pro diplomovou práci je ISO 690. Jména časopisů lze uvádět zkráceně, ale jen
%%% v kodifikované podobě. Všechny použité zdroje a prameny musí být řádně citovány.

\def\bibname{Bibliography}
\begin{thebibliography}{99}
\addcontentsline{toc}{chapter}{\bibname}

\bibitem{etoile_article}
  David Chisnall,
  \emph{A Modern Objective-C Runtime}.
  2009.
  \url{http://www.jot.fm/issues/issue_2009_01/article4/index.html}
  
\bibitem{openstep_wiki}
  \emph{OpenStep}.
  \url{http://en.wikipedia.org/wiki/OpenStep}
  
\bibitem{rhapsody_wiki}
  \emph{Rhapsody (operating system)}.
  \url{http://en.wikipedia.org/wiki/Rhapsody_(operating_system)}
  
\bibitem{mountain_lion_enhancements}
  \emph{OS X 10.8 run-time enhancements}.
  \url{http://cocoaheads.tumblr.com/post/17719985728/10-8-objective-c-enhancements}
  
\bibitem{smalltalk_class_diagram}
  Andrew P. Black St\'{e}phane Ducasse, Oscar Nierstrasz Damien Pollet, Damien Cassou and Marcus Denker,
  \emph{Smalltalk Class Diagram}.
  \url{http://pharo.gforge.inria.fr/PBE1/PBE1.html}
  
\bibitem{apple_open_source}
  \emph{Apple Open Source}.
  \url{http://www.opensource.apple.com}  

\bibitem{safari_windows}
  Jong Am,
  \emph{Safari 4 beta for Windows}.
  \url{http://jongampark.wordpress.com/2009/02/24/safari-4-beta-for-windows/}  
 
\bibitem{apple_malloc_man_page}
  \emph{malloc\_zone\_malloc(3) OS X Developer Tools Manual Page}.
  \url{http://developer.apple.com/library/Mac/\#documentation/Darwin/Reference/ManPages/man3/malloc_zone_malloc.3.html}

\bibitem{class_dump}
  Steve Mugard,
  \emph{class-dump}.
  \url{http://stevenygard.com/projects/class-dump/} 
  
\bibitem{toll_free_bridging_internals}
  Mike Ash,
  \emph{Friday Q&A 2010-01-22: Toll Free Bridging Internals}.
  \url{http://www.mikeash.com/pyblog/friday-qa-2010-01-22-toll-free-bridging-internals.html} 

\bibitem{thread_local_storage}
  Andy Monitzer,
  \emph{Thread-Local Storage}.
  \url{http://gcc.gnu.org/onlinedocs/gcc-4.3.2/gcc/Thread_002dLocal.html} 
 
\bibitem{java_bridge}
  Andy Monitzer,
  \emph{The Java Bridge}.
  \url{http://cocoadevcentral.com/articles/000024.php} 
  
\bibitem{associated_objects}
  \emph{}.
  \url{https://developer.apple.com/library/mac/\#documentation/Cocoa/Reference/ObjCRuntimeRef/Reference/reference.html#//apple_ref/c/func/objc_setAssociatedObject} 

\end{thebibliography}


%%% Figures used in the work.
\addcontentsline{toc}{chapter}{\listfigurename}
\listoffigures{}

%%% Použité zkratky v diplomové práci, existují-li, včetně jejich vysvětlení.
\chapwithtoc{List of Abbreviations}

%%% Přílohy k diplomové práci, existují-li (různé dodatky jako výpisy programů,
%%% diagramy apod.). Každá příloha musí být alespoň jednou odkazována z vlastního
%%% textu práce. Přílohy se číslují.
\chapwithtoc{Attachments}

\openright
\end{document}
